% Bu dosya, Bolum 13'u zenginlestirir (Yeni chapter acmaz)
\section{OSINT Teknikleri ve Kurumsal Farkındalık Programı}

\subsection{OSINT Pratikleri}
\begin{itemize}
  \item \textbf{E-posta ve Alan Adı Ayakizi:} theHarvester, amass, crt.sh
  \item \textbf{Sosyal Medya Ayakizi:} Kurumsal sızıntılar, unvan/rol değişimleri
  \item \textbf{Döküman Metaverisi:} exiftool ile sızabilecek dahili bilgiler
\end{itemize}

\Not{OSINT çalışmaları yasal ve etik çerçevede yapılmalıdır. Kurum içi tatbikatlarda yazılı onay ve kapsam tanımı şarttır.}

\subsection{Kırmızı Takım Sosyal Mühendislik Playbook'u}
\begin{enumerate}
  \item \textbf{Keşif:} Personel listeleri, tedarikçi ilişkileri, etkinlik takvimi
  \item \textbf{Kandırma Senaryosu:} Otorite, aciliyet ve karşılık verme prensiplerini içeren bağlamlı kurgular
  \item \textbf{Kanal Seçimi:} E-posta, telefon (vishing), sahte toplantı daveti, fiziksel erişim
  \item \textbf{Kontrol Noktaları:} Güvenli kelime, escalation akışı, durdurma koşulları
  \item \textbf{Ölçüm:} Tıklama oranı, kimlik bilgisi ifşası, raporlama oranı, düzeltici aksiyon
\end{enumerate}

\Uyari{Gerçek sistemlere yetkisiz erişim denemesi ve kişisel verilerin izinsiz işlenmesi suçtur. Tatbikatlar denetimli ortamda yürütülmelidir.}

\subsection{Farkındalık Programı Çerçevesi}
\begin{itemize}
  \item \textbf{İçerik:} Kuruma özel oltalama örnekleri, gerçek vakalar, rol tabanlı modüller
  \item \textbf{Sıklık:} Aylık mikro içerikler + yıllık kapsamlı eğitim
  \item \textbf{Ölçüm:} Eğitim tamamlama, test puanları, oltalama tatbikatı metrikleri
  \item \textbf{Süreklilik:} Yeni tehditlere göre güncellenen içerik ve kısa “anlık uyarılar”
\end{itemize}

\Ipucu{Farkındalık eğitimlerini HR/L\&D süreçlerine entegre edin ve yöneticilere özel özet panolar sağlayın.}
